\documentclass[11pt]{article}

\usepackage[margin=1.1in]{geometry}
\usepackage{amsmath,amssymb,amsthm}
\usepackage{mathtools}
\usepackage{stmaryrd}
\usepackage{microtype}
\usepackage[T1]{fontenc}
\usepackage{palatino}
\usepackage{enumitem}
\usepackage[hidelinks]{hyperref}

\theoremstyle{definition}
\newtheorem{definition}{Definition}
\newtheorem{principle}{Principle}
\theoremstyle{plain}
\newtheorem{proposition}{Proposition}
\newtheorem{conjecture}{Conjecture}
\newtheorem{question}{Question}
\theoremstyle{remark}
\newtheorem{remark}{Remark}

% process-calculus notation
\newcommand{\deref}{{*}}
\newcommand{\inn}[3]{\mathrm{for}(#1 \leftarrow #2)\,#3}
\newcommand{\out}[2]{#1\,!\,(#2)}
\newcommand{\nil}{\mathbf{0}}
\newcommand{\rcong}{\equiv}
\newcommand{\bisim}{\sim}
\newcommand{\rbisim}{\sim_{r}}
\newcommand{\reduces}{\longrightarrow}
\newcommand{\sem}[1]{\llbracket #1 \rrbracket}
\newcommand{\bagunion}{\uplus}
\newcommand{\Name}{\mathrm{Name}}
\newcommand{\Proc}{\mathrm{Proc}}
\newcommand{\dgr}{\mathbf{d}}
\newcommand{\Cut}{\mid}
\newcommand{\GSLT}{\mathsf{GSLT}}
\newcommand{\Hyp}{\mathsf{Hyp}}

\title{\bfseries Where Does Agency Live?\\[3pt]
\large Nondeterminism as a computational resource,\\
and agency across the computational and biological axes}

\author{Lucius Gregory Meredith\\
\small CEO, F1R3FLY.io, and Chief Science Officer, F1R3FLY Industries\\
\small \texttt{lgreg.meredith@gmail.com}}

\date{2026}

\begin{document}
\maketitle

\begin{abstract}
The relationship between agency and determinism is an old question; this note
brings a new instrument to it. We plot the phenomena on two axes ---
computation and biology --- and use each to calibrate the other. The
analytical hinge is a single reframing: in a model of computation where
interaction is primitive, the nondeterminism of an interaction is not missing
information but a \emph{resource}, namely a choice. Made precise, that choice
is a choice function for a definite family, so the mathematics of the axiom of
choice acquires a computational meaning, and the strength of the choice
principle one spends becomes a coordinate in a hierarchy of computational
power --- a stalk in a \emph{fibration of hypercomputation} over the category
of models of computation and their bisimulation-preserving morphisms. The
mobile process calculi sit in this category as the objects built expressly for
agency: they take all computation to arise from interaction among agents, and
they are \emph{fundamentally non-confluent}, which is exactly what is needed to
speak of races --- two sodium atoms for one chlorine, two sperm for one egg,
two predators for one prey, two buyers for the last ticket. RSpace, an
abstraction that lets these calculi run in production, is where the choice
implicit in the nondeterminism becomes visible. We then turn the instrument on
agency itself: once choice is localised, the question of where to draw an agent's
boundary gets a fresh answer. Eric Smith's reading of the biosphere is our point
of entry --- he argues that we may have to be individual-organism-blind and treat
the whole biosphere as the agent --- but we take from him only the lesson that the
boundary is negotiable, not his placement of it: his interior, full of predation
and natural selection, is anything but determinate. We redraw the boundary by a
different and finer criterion --- confluence up to resource bisimilarity ---
around the regions that carry no observable choice, branch as they may, in the way
the $\lambda$-calculus is determinate. Such regions, macro-components, turn out to
be small, local islands in a dense sea of races; Smith's thermodynamic whole and
our determinacy islands cross rather than nest. The
fibration, finally, is a map and not the territory: the actual universe of
computational phenomena may resemble the physical one, with scales and capacities
wildly juxtaposed --- femtosecond molecular dynamics beside minute-scale gene
expression in a single cell, halting oracles beside $\lambda$-terms at a single
cut. That juxtaposition is the flat ontology, and biology is our reason to take
it seriously. If choice in nondeterministic interaction is how hypercomputation
would show up physically, it would be correspondingly hard to see --- but Bell's
correlated choice, formalised and tested in the laboratory, is the precedent that
says it need not be invisible, and that the mathematics of choice and a
computationally informed physics may have much to give each other.
\end{abstract}

\section{An old question, two axes}

Is agency compatible with a determined world, and if there is room for agency,
where exactly is the room? The question is old and the usual moves are
well-worn. We want to bring a different instrument: the theory of computation,
sharpened by a particular family of models in which interaction is the primitive
and nondeterminism is not noise but structure --- and, in parallel, the biology
of living matter, where agency is not a philosopher's posit but the daily
business of cells, organisms, and ecosystems. We treat these as two axes on which
to plot the same phenomena, and the discipline we hold ourselves to is
bidirectional: every claim should read in both directions. Computation should
inform our picture of biological agency, and biology should inform our picture of
what computation is. The flat ontology we arrive at is motivated precisely by
that second direction --- the living world shows us, concretely, the
juxtaposition of scales and capacities that the computational picture, left to
its own neat hierarchies, would not have predicted.

The fresh move is small to state and large in consequence. In a model of
computation built on interaction, when two agents meet and could combine in more
than one way, the selection of \emph{which} way is a genuine degree of freedom.
The standard reflex is to treat that degree of freedom as something to be averaged
away or refined out --- nondeterminism as ignorance. We treat it instead as a
\emph{resource}: a choice, in the precise mathematical sense, whose supply has a
cost and whose strength can be measured. Once nondeterminism is a resource, two
things follow that organise the whole essay. First, the discourse of the axiom of
choice --- usually filed under foundations of mathematics --- becomes
computationally load-bearing: different ways of supplying choice import different
amounts of computational power. Second, the question of agency relocates. Agency
lives where the choice is open, and the boundary of an agent is best drawn around
the cuts where its choices are genuinely live. Not every place two agents meet is
such a place; much of what looks like interaction is determinate traffic with no
choice to be had. Finding where the choice actually is turns out to be the same
as finding where agency actually is.

\section{A landscape of computational models}\label{sec:landscape}

To make ``a coordinate in a hierarchy of computational power'' precise we need a
space of models, not a single model. Fix the following landscape, briskly.

\paragraph{Objects: models of computation as graph-structured lambda theories.}
A wide range of computational formalisms --- the $\lambda$-calculus, term-rewriting
systems, Petri nets, the process calculi --- can be presented uniformly as
\emph{graph-structured lambda theories} (GSLTs): a syntax of terms with a binding
discipline, a notion of context, and a reduction or interaction relation, all
organised so that the term structure is a graph rather than merely a tree
\cite{oslf}. The details of the presentation are not needed here. What we need is
that each such model carries a canonical behavioural equivalence ---
\emph{bisimilarity} --- generated, in the GSLT discipline, by labelling transitions
not with primitive actions but with the minimal contexts that enable them, after
Leifer and Milner \cite{leifermilner,sewell}. Two computations are bisimilar when
no context can tell them apart; the observer-specified logic functor (OSLF) of
\cite{oslf} generates, from the same context labels, the modal logic whose formulae
they satisfy, so behavioural equality and logical indistinguishability coincide.

\paragraph{Morphisms: maps that preserve bisimulation.}
Let $\GSLT$ be the category whose objects are these models and whose morphisms are
the maps that preserve bisimilarity --- translations of one model into another that
respect when two computations are behaviourally the same. This is the right notion
of morphism for our purpose, because the thing we will fiber over each model is its
\emph{observable behaviour}, and behaviour is exactly what bisimilarity tracks. A
bisimulation-preserving morphism carries a model's behavioural fabric, and hence its
hierarchy of behaviours, into another's.

\paragraph{The fibration of hypercomputation.}
Over the single Turing-complete model, Turing already gave us a tower: relativise
computation to an oracle, take the halting problem of the relativised machine as a
new oracle, iterate, and index the iteration by the ordinals \cite{turing1939} ---
the arithmetical hierarchy and the Turing degrees, the classical apparatus of
hypercomputation \cite{soare,copeland}. The instrument we want is the
\emph{generalisation of this tower to a tower per model}. Over each object $G$ of
$\GSLT$ stands not a single computational power but a whole hierarchy ---
$G$ relativised to its own oracles, the hypercomputation tower \emph{specific to that
model of computation}. These towers assemble into a fibration
\[
p \colon \Hyp \longrightarrow \GSLT,
\]
whose fiber (stalk) over $G$ is the hypercomputation hierarchy native to $G$, and
whose vertical direction is added computational power. Reindexing along a
bisimulation-preserving morphism $G' \to G$ --- the cartesian direction --- carries
a high-power computation to its shadow in a weaker frame, which is exactly the act
of describing a strong system in a weaker vocabulary. We will lean on this picture
twice: first to give ``choice as a resource'' a vertical coordinate, and later, when
we ask whether the fibration is the territory or only a chart of it.

\section{Interaction, and why non-confluence is the point}\label{sec:interaction}

Among the objects of $\GSLT$, one class is built for our question. The mobile
process calculi (MPCs) --- Milner's $\pi$-calculus and, in the form we use, the
reflective higher-order (rho) calculus \cite{rho} --- take a definite philosophical
stance: \emph{all computation arises from interaction among agents}
\cite{milnerCC}. There is no lone machine grinding a tape; there are processes that
meet, exchange, and continue, and computation is what happens at their meetings. We
write the meeting as a parallel composition, the \emph{cut}:
\[
P \Cut Q,
\]
and the single base reaction is communication across it,
\[
\inn{y}{x}{P} \Cut \out{z}{Q} \reduces P\{@Q/y\} \qquad (x \equiv_N z),
\]
an input on a channel meeting a matching output. Agents flank every cut; the cut is
where they act on one another.

The feature that makes the MPCs the right tool --- and that distinguishes them
sharply from the $\lambda$-calculus --- is that they are \emph{fundamentally
non-confluent}. When several outputs and several inputs collect on the same channel,
which output meets which input is not determined by the terms; it is a real choice,
and different resolutions lead to genuinely different, non-reconvergent futures.
This is usually presented as a complication. We present it as the entire point,
because non-confluence is the computational signature of a \emph{race}, and races
are how agency shows up in the world. Consider, across four scales:

\begin{itemize}[leftmargin=1.4em]
\item two sodium atoms competing to donate their electron to one chlorine;
\item two sperm racing to fertilise one egg;
\item two predators converging on a single prey;
\item two buyers reaching for the last concert ticket.
\end{itemize}

\noindent
Each is, structurally, one consume-slot (the chlorine, the egg, the prey, the
ticket) contended by several produce-candidates, resolved by a pairing that the
situation does not fix, with the unmatched candidates left over. Each outcome
\emph{matters}: this sodium and not that one is now bound; this sperm and not that
one founds the lineage. A confluent model cannot speak of these, not because it
lacks the vocabulary but because confluence is precisely the claim that the choice
does not matter --- that all resolutions reconverge. The $\lambda$-calculus, the
great confluent model, is the stakes-free corner of the landscape: its many
reduction orders are real choices whose outcomes are identified by Church--Rosser,
so the choice is free and invisible. Reality is mostly not like that. To talk about
agency one needs a model in which the choice can matter, and the MPCs are the
objects of $\GSLT$ in which it does. This is the first crossing of our two axes:
biology (the race for the egg) tells us which computational models are adequate to
agency, and the answer is the non-confluent ones.

\paragraph{The rho calculus, specifically.}
We use the rho calculus because its names are quoted processes --- $@P$ turns a
process into a name, $\deref x$ turns a name back into a process, and
$\deref @P \equiv P$ --- so there are no primitive names and no name-restriction
operator; reflection furnishes recursion with no primitive replication \cite{rho}.
Nothing below turns on the set-theoretic models of \cite{qcs,knot}; we draw only on
the shape of the calculus. The reader may keep the modern surface syntax in view:
$@P$ for quote, $\deref x$ for dereference, $\inn{y}{x}{P}$ for input, $\out{x}{Q}$
for output, $\Cut$ for the cut, $\rcong$ for structural congruence, $\bisim$ for
context bisimilarity (the behavioural, congruence-by-construction equivalence) and
$\rbisim$ for resource bisimilarity (the finer equivalence that counts parallel
copies).

\section{RSpace: the bridge, and where the choice is}\label{sec:rspace}

The MPCs are admirable for reasoning and awkward to run. RSpace is the abstraction
that closes the gap: it presents the parallel fragment of the rho calculus as a
store --- a hashtable from channels to polarity-homogeneous bags, either of data
(produces) or of continuations (consumes) --- and interprets the cut as keyed
multiset union, so the denotation $\sem{-}$ is a monoid homomorphism and the
structural laws of $\Cut$ become literal equalities of tables \cite{qcs}. This is
not a toy: it is how the calculus is compiled into a high-performance, persistent,
content-addressed store in production. We invoke it here for one reason: it makes
the choice of \S\ref{sec:interaction} \emph{visible and operational}.

When a channel carries both a bag of consumes and a bag of produces, communication
chooses an underlying list of each and \emph{zips} them --- matching pairs react,
unmatched payloads remain as surplus at the channel. The denotation stays a
function; the zip is the separate dynamics; and the choice of which list orderings
to zip is the calculus's nondeterminism, now pinned to a key. The four races of
\S\ref{sec:interaction} are this forgiving zip with surplus, drawn from life:
several candidates, one slot, a pairing the data does not fix, leftovers remaining.
RSpace's contribution to the present inquiry is to show that the nondeterminism is
not a haze over the system but a specific, locatable act of selection --- and an act
of selection is a choice, in the sense the next section makes exact.

\begin{remark}[A motivating example, not a mechanism]
We are not proposing that sodium, sperm, or ticket-buyers \emph{run} RSpace, or that
its store is a physical mechanism. The claim is structural: the abstract shape ---
contended slot, nondeterministic pairing, surplus --- is shared, and RSpace is where
that shape, and the choice inside it, is displayed cleanly enough to analyse. The
value of the MPC and RSpace vocabulary is that it \emph{names the categories};
whether any given physical system instantiates them is a separate, empirical
question we do not prejudge.
\end{remark}

\section{Choice as a computational resource}\label{sec:choice}

Now make the choice exact and connect it to the stalks of \S\ref{sec:landscape}.
Index by the channels carrying mixed polarity (or, for a run, by the reaction
events). At each index the fibre is the set of admissible pairings ---
inhabited precisely when a redex is present. A completed schedule is a section of
this family: an element of the product of the fibres. The assertion that a section
exists is an instance of the axiom of choice, and the scheduler that realises it is,
literally, the choice function for the pairing family. Nondeterminism resolved
\emph{is} choice exercised. This is the practical reframing of $\mathrm{AC}$: not a
remote axiom about arbitrary sets but the principle that supplies a scheduler, and
its various restrictions and strengthenings are different schedulers with different
powers.

The strength required is graded, and the grading is the one the source calculus's
risk ledger already keeps \cite{qcs}.

\begin{itemize}[leftmargin=1.4em]
\item \emph{Finite contention, finite run.} The index set is finite and the product
is inhabited with no choice axiom at all --- constructively. More: a bag is a
parallel composition of finitely many prefixes, so it arrives \emph{presented} as a
list-up-to-permutation. The term hands you the enumeration; choosing a linearisation
is free. In the finite world choice is not non-constructive; the term is its own
choice structure.

\item \emph{One $\omega$.} Reflection (or non-termination) installs an
$\omega$-fold contention, or an $\omega$-sequence of events whose pairings depend on
the residuals of the last. The static case needs countable choice; the dependent case
needs Dependent Choice. This single $\omega$ is the one the ledger spends; here it
buys exactly the step from no-choice to countable/dependent choice.

\item \emph{Class size.} Full $\mathrm{AC}$, and only here.
\end{itemize}

\noindent
The decisive connection is vertical. A scheduler that supplies only finite,
demonic choice keeps the computation at the base of its hypercomputation stalk: it
imports no power the model did not already have. A scheduler that can survey an
infinite contention and pick a globally coherent matching --- an angelic or
oracle-endowed scheduler --- imports computational power, and that power is altitude
in the fiber over the model. The cleanest way to say it: \emph{choice is the vertical
coordinate}. Where you sit in the hypercomputation stalk of \S\ref{sec:landscape} is
set by how strong a choice principle your scheduler is allowed to spend. The Weihrauch
lattice \cite{weihrauch}, whose objects are choice principles ordered by uniform
reducibility and composed by ``do one, then the other,'' is the natural home for this
gradation --- and, tellingly for what follows, it is a \emph{lattice}, not a linear
chain. This is the second crossing of our axes, now internal to computation: the
foundational mathematics of choice turns out to measure the same thing as the
theory of relative computability, because both are measuring the resource a race
consumes when it is resolved.

\section{Choice and the drawing of boundaries}\label{sec:boundaries}

We can now turn the instrument on agency. If nondeterminism is localised at cuts and
is a resource, then the natural place to look for an agent's boundary is where its
choices are live. Make two notions precise.

\begin{definition}[Cut and aperture]\label{def:aperture}
A \emph{cut} is a parallel composition co-locating a consume and a produce --- the
locus where two agents interact. A cut is an \emph{aperture} when the resolution of
its pairing nondeterminism is \emph{observable}: distinct admissible pairings lead to
residuals in distinct behavioural classes (distinct $\rbisim$-classes at the resource
layer, or distinct $\bisim$-classes after forgetting multiplicity). A cut that is not
an aperture carries no observable choice: either there is a unique match, or the
several matches reconverge.
\end{definition}

\begin{principle}[Not every cut is an aperture]\label{prin:main}
Agents flank every cut, but only apertures leave the interaction's nondeterminism
open. And only at an aperture can an agent of higher computational capacity influence
the behaviour of a lower one --- by biasing the live choice. Where a cut is not an
aperture there is no choice to bias, and the stronger agent's surplus power finds no
purchase. Agency, in the sense of room-to-do-otherwise, lives at the apertures.
\end{principle}

\noindent
This is already a fresh answer to the old question. The apertures are exactly the
races; the determinate traffic between them is not where agency is decided, however
much computation flows through it. And it locates interference precisely: a
high-capacity agent reaches into a low-capacity one not everywhere they touch but
only at the cuts where the low-capacity one has a live, observable choice that the
high-capacity one can see how to resolve. The asymmetry of capacity --- one agent
several halting-oracles above another --- becomes operative only through an aperture.

Two corrections to a tempting oversimplification, both of which we will need. First,
the aperture is a property of the \emph{cut}, not of either agent; nondeterminism is
not preserved by $\Cut$, so a determinate whole can have wildly underdetermined
internal cuts, and conversely. ``An agent's nondeterminism budget is its interface''
is therefore not quite right: the budget belongs to the cut, and which cut one calls
the boundary is a modelling decision. Second, because of this, a seam where choice is
genuinely open may sit far from the boundary an observer would draw, and its effects
may be felt at cuts the observer counts as internal. We turn to a setting where this
plays out concretely.

\section{The biological axis: Smith redraws the boundary}\label{sec:smith}

Smith and Morowitz read terrestrial life as a planetary process and ask after its
interface to the rest of the universe \cite{smith}. On their account the agents that
matter at that interface are the ones that consume energy from \emph{outside} the
biosphere --- sunlight, the chemical disequilibria of hydrothermal vents --- while a
vast interior of agents consume \emph{one another}, trading energy already captured.
The energy-eaters are the actual boundary between (terrestrial) life and the cosmos,
and Smith presses the point we want: that we may have to be
\emph{individual-organism-blind} and treat the entire biosphere as the agent. The
naive ontology places the agency boundary at the organism --- this bee, that bird ---
and Smith argues that on a thermodynamic criterion the boundary belongs somewhere else
entirely, around the whole. For him the biosphere is the macro-component, unified by
its single interface to the extra-biospheric energy source.

This boundary-redrawing is exactly the move we need, and it is the only thing we take
from the analogy. We do \emph{not} take Smith's interior to be determinate. It plainly
is not. Schematically,
\[
\textsf{Biosphere} \;\rcong\; \textsf{EnergyEaters} \Cut \textsf{AgentEaters},
\qquad
\textsf{Universe} \;\rcong\; \textsf{Biosphere} \Cut \textsf{RestOfUniverse},
\]
the seam to the cosmos being $\textsf{Biosphere} \Cut \textsf{RestOfUniverse}$ and
\textsf{EnergyEaters} the sub-term whose business is that cut --- where outside energy,
untyped to any particular molecule (the sun does not address its photons to a chosen
chlorophyll), enters as an underdetermined produce, a genuine aperture. But
\textsf{AgentEaters} is \emph{full} of meaningful races, exactly the ones
\S\ref{sec:interaction} put on the table: predation contended (two predators, one
prey), competition for scarce resource, and, most decisively, \emph{natural selection
and speciation}, which are non-confluent almost by definition --- which variant fixes,
which lineage splits, is precisely an outcome where the choice matters and the
histories do not reconverge. Smith's biosphere-agent is therefore \emph{not} a
macro-component in our sense. Its interior runs the deepest races biology has.

So the relationship between Smith's factorisation and ours is not that we fill in his
interior; it is that we propose a \emph{different criterion} for the same kind of move.

\begin{remark}[Two boundary-redrawings; ours strictly finer]\label{rmk:refine}
Smith redraws the naive per-organism boundary \emph{outward}, to one biosphere-sized
agent, on a thermodynamic criterion: what interfaces with the external energy source.
We redraw it by a \emph{different} criterion --- confluence up to resource bisimilarity
(\S\ref{sec:macro}) --- which lands the boundary neither at the organism nor at the
whole, but around whatever subsystems carry no observable choice. This is strictly more
refined: it asks not ``what touches the outside?'' but ``where is there no live race?''
The two need not nest. Smith's macro-component contains selection, so it is not one of
ours; our macro-components, as the next section argues, are small and local, so the
biosphere is not one of theirs. Smith's whole and our islands are different objects,
unified by different things --- thermodynamics versus determinacy --- and the honest
statement is that the factorisations \emph{cross} rather than refine one within the
other. What we credit Smith with, and it is the load-bearing credit, is the
demonstration that the naive placement of the agency boundary is \emph{negotiable}:
``agent'' need not mean ``organism.'' Once that is conceded, the question is which
criterion should fix the boundary, and we answer it differently.
\end{remark}

\noindent
This still gives the seam-versus-boundary phenomenon of \S\ref{sec:boundaries} its
biological face, but more carefully than a single clean interface would. The energy cut
is an aperture; its bias propagates into a teeming interior that has plenty of its own
apertures; and the genuine computational question --- whether such a bias is, at any
given depth, transmitted through a determinate region or absorbed into a fresh race ---
is answered locally, region by region, not once for the whole. To answer it at all we
need the determinate regions named precisely, and they are smaller than Smith's whole.

\section{Back to computation: confluent macro-components}\label{sec:macro}

We want a name for the determinate regions \S\ref{sec:smith} promised: subsystems all
of whose internal cuts are sealed, with apertures only at the boundary. These are not
Smith's biosphere --- which we just conceded is race-riddled --- but the genuinely
choice-free pockets, wherever they lie. The first guess is that such a region is sealed
by \emph{rigidity} --- a unique match at every internal cut, so the zip is a function
and there is no race. Syntactically that is \emph{linearity}: each internal channel used
once per polarity, each consume discharged by exactly one produce, the boundary the
residue of the contended channels. But linearity is too strong. What we actually require
is only that the interior carry no \emph{observable} choice, and the exact statement of
that is confluence.

\begin{definition}[Macro-component]\label{def:macro}
A subsystem is \emph{rigid} if its internal reaction is confluent up to resource
bisimilarity: every internal race reconverges, all maximal internal reaction
sequences reaching $\rbisim$-equal normal forms. A \emph{macro-component} is a maximal
rigid subsystem; its internal cuts are non-apertures and its boundary cuts --- where
internal confluence fails --- are its apertures. The locus of confluence-failure is the
\emph{rigidity frontier}.
\end{definition}

\noindent
Confluence permits the branching linearity forbids. A bee may have several
internally admissible pairings; if all flow to $\rbisim$-equal normal forms it is no
site of observable choice, which is all the factorisation needs. Metabolism is full of
internal branching --- redundant pathways, order-independent cascades --- that
converges on the same steady state: confluent, not linear. The slogan is therefore
\emph{the interior is $\lambda$-like, not linear-like}. The $\lambda$-calculus is the
witness: operationally awash in branching, one reduction order against another, and
famously determinate because Church--Rosser makes the branching fictional --- one
normal form, reached however you like. A \emph{closed} $\lambda$-term is a sealed
macro-component at every scale, nothing to interfere with because no channel is
contended; an \emph{open} term is the agent with apertures, its free variables the
boundary channels where a context can substitute and change the result. In the MPC
vocabulary the free variable is the unmatched consume, the dangling polarity --- the
port where the world reaches in.

The picture this forces is not one determinate interior behind one skin. If selection,
predation, and competition are non-confluent and pervasive --- and \S\ref{sec:smith}
grants that they are --- then confluent macro-components are \emph{small and local}: a
tightly coupled metabolic or developmental subsystem whose internal branching genuinely
reconverges (redundant enzymatic pathways, buffered developmental cascades), bounded by
an aperture wherever a real race begins. The biosphere is then a vast field of small
confluent islands in a dense sea of races, rather than a single interior. This is less
tidy than a clean interface and it is the true shape: \emph{apertures are everywhere},
the determinate regions are the exception, and ``where is the agent?'' has many local
answers rather than one global one. An ecosystem may be the right scale for some such
island --- but only the parts of its food web that reconverge up to resource
bisimilarity, with its contended predation and its competition for influx counted as
apertures \emph{at or within} it, not as interior. The ecological races are not
quarantined behind the boundary; they are the boundary, threaded through. Whether
\emph{any} biosphere-scale structure is confluent, or whether confluence is strictly a
small-scale phenomenon with all large structure race-dominated, we leave open
(\S\ref{sec:open}); if the latter, the macro-component is real but always local, and
Smith's whole and our islands simply do not nest.

This is where two senses of ``agent'' come apart, and the distinction is the conceptual
payoff.

\begin{remark}[Two senses of ``agent'']\label{rmk:two}
A confluent island's internal parts may be sharply individuated --- distinct enzymes,
distinct cell types, distinct roles --- without being \emph{sites of choice}. Their
distinctness is type-distinctness, not aperture-distinctness; whatever internal branching
they run reconverges, so no internal cut is a live race. Discrimination among agents and
the location of nondeterminism therefore come apart: an agent in the \emph{discrimination}
sense (an individuated role) need not be an agent in the \emph{choice} sense (a cut with a
live race), and vice versa --- a featureless region can sit astride a fierce race. A finer
description sees \emph{more} individuated agents but not necessarily more apertures; the
apertures are wherever the races are, which is a separate question from where the roles
are. Naive ontology conflates the two, placing ``the agent'' at the individuated role.
Both Smith's move and ours are refusals of that conflation --- his relocating the boundary
by thermodynamics, ours by determinacy.
\end{remark}

\noindent
With the determinate regions named, the cuts that matter for agency are the apertures
between and within them --- where the nondeterminism of \S\ref{sec:interaction} is real.
Three frontiers coincide there, which is the structural dividend:

\begin{proposition}[One frontier, three readings]\label{prop:frontier}
For a macro-component, the rigidity frontier (where internal confluence up to
$\rbisim$ fails) coincides with the \emph{detectable-interference frontier} (where a
biased pairing yields a different $\rbisim$-class, so an external write leaves a mark)
and with the \emph{confluence-obligation frontier} (where the eager-confluence
obligation of the path semantics \cite{paths} is non-trivial). The place where the
interior can no longer hide its schedule is the place where interference becomes
observable, which is the place where the formal obligation bites.
\end{proposition}

\noindent
A consequence worth isolating: a confluent interior carries no choice resource --- no
observable nondeterminism, hence no Weihrauch content, hence (\S\ref{sec:choice}) no
altitude in the hypercomputation stalk --- however baroque its branching. So the
resource degree of a macro-component is carried entirely by its boundary,
$\deg(\text{macro-component}) = \deg(\text{boundary})$. Capacity, like choice, lives at
the apertures.

\section{The flat ontology}\label{sec:flat}

The fibration of \S\ref{sec:landscape} is a map, and we should not mistake it for the
territory. It is a way to \emph{understand} computational phenomena --- to sort agents
by capacity and read interference as reaching down a tower. But the actual universe of
computational phenomena may be far less tidy, and biology is the reason to suspect so.
In a single living cell, fast molecular dynamics --- vibrational and electronic
transitions down to the femtosecond --- coexist with transcription and translation
that unfold over seconds to minutes: many orders of magnitude apart, intimately
coupled, in one small volume. The cell does not run its fast and slow processes in
separate strata that communicate through a clean interface; it juxtaposes them, and the
juxtaposition is the point. The physical universe at large is the same --- scales and
scopes wildly interleaved rather than cleanly layered.

We propose to take computation the same way. In the \emph{flat ontology}, the universe
of computational phenomena is the total space of the fibration \emph{without privileging
the projection} --- a jumble of agents of wildly different computational capacity,
several halting-oracles up beside fully deterministic $\lambda$-like agents,
interacting directly at cuts. The tower is not the structure of the world; it is a
cleavage, a chosen gauge for describing a frame-independent jumble, and the
``stratification'' is an artefact of which morphism we chose to project along. The
islands-in-a-sea-of-races picture of \S\ref{sec:macro} is the same claim seen from the
agency side: determinate regions are local and sparse, apertures pervade, and there is no
clean global stratum at which all the choice has been pushed to one rim. Read this in
both directions, as promised. Computation, taught by biology, stops expecting a neat
hierarchy and starts expecting juxtaposition. Biology, taught by computation, gains a
vocabulary for what the juxtaposition costs and where it can bite: a high-capacity agent
can influence a low-capacity one, but --- Principle~\ref{prin:main} --- only at an
aperture, only where the lower agent has a live choice the higher one can resolve.

What keeps such a flat universe coherent, rather than collapsing to the top capacity?
Two disciplines, both already in the MPC vocabulary. \emph{No implicit interaction}: the
store never reacts with itself; a reaction is triggered only by a cut that actually
co-locates opposite polarities \cite{paths}. Drop it and a single oracle on a universal
channel homogenises the field upward; keep it and interference must be \emph{cut-triggered}.
\emph{Locality and cost}: a strong agent can write only where it shares a channel ---
and, in the path reading, only into the subspace below a prefix it holds --- and a cost
discipline makes its reach finite per unit resource. Locality plus cost is what lets
agents of disparate capacity share a world without the strongest simply absorbing the
rest. These are not bookkeeping; they are the consistency conditions for a flat ontology
of unequal agents --- and they are exactly the conditions under which a small confluent
island can persist amid the surrounding races without being dissolved by the strongest
agent that happens to share one of its channels.

\section{The hinge: is determinism stable under interference?}\label{sec:crux}

One question decides whether the picture is static or dynamic, and it is the question a
formalisation should resolve first. Everything above treats ``determinate interior,
nondeterministic boundary'' as a property a system \emph{has}. Whether it \emph{keeps}
the property when a boundary aperture fires and the decision propagates inward is not
obvious --- because confluence up to resource bisimilarity, unlike linearity, is not
compositional, and because a deep interaction in the path semantics manufactures a fresh,
deeper interaction whose confluence is a new obligation \cite{paths}.

\begin{question}[Stability of rigidity]\label{q:crux}
Is confluence up to resource bisimilarity preserved as a boundary decision cascades into
the interior? Equivalently: when a macro-component's aperture fires, does its interior
stay rigid?
\end{question}

\noindent
The two answers are both meaningful, and the choice between them is what biology and
computation are jointly asking. If rigidity is stable, the factorisation is robust: a
boundary race fires, the decision descends as a \emph{forced} cascade through a confluent
island that has no opinion of its own, and ``determinate interior / nondeterministic
boundary'' is a genuine structure theorem holding along the dynamics --- a buffered
developmental cascade absorbing an upstream choice into one canalised outcome. If rigidity
is \emph{not} stable, then boundary interference can crack open a previously determinate
island, turning a confluent agent into a branching one by descending into it. That is not
merely a defect;
it is plausibly the formal content of \emph{perturbation} --- a high-capacity agent not
only steering a low-capacity one within its existing freedom but \emph{manufacturing} new
freedom in it and then steering that. In the ontological register: the difference between a
universe whose agents have fixed capacities and one in which capacity itself is something
agents can do to one another. We do not know which holds, and the same confluence lemma
settles it on both axes at once.

\section{Open problems}\label{sec:open}

We collect the load-bearing uncertainties, each a structural claim one can try to settle.
\emph{(i)} Make precise the family-of-pairings whose section is the schedule and prove the
$\mathrm{AC}$-grading of \S\ref{sec:choice}, identifying which schedulers are confluent up
to $\rbisim$ (constructive, no imported oracle) and which import altitude in the stalk.
\emph{(ii)} Exhibit the fibration $p\colon\Hyp\to\GSLT$ rigorously and test whether angelic
and demonic resolution are the Lawvere adjoints $\exists \dashv (-)^{*} \dashv \forall$ to
reindexing --- the decisive sub-question being whether the GSLT fibres carry the limits and
colimits both adjoints require. \emph{(iii)} Show ``maximal rigid subsystem''
(Definition~\ref{def:macro}) is well defined and that the three frontiers of
Proposition~\ref{prop:frontier} coincide, and establish that a confluent interior is
resource-degree-transparent. \emph{(iv)} Resolve Question~\ref{q:crux} --- the stability of
rigidity under descent --- proving the structure theorem if it holds, or characterising
perturbation as a property of the sub-system if it fails. \emph{(v)} Settle the scale
question of \S\ref{sec:macro}: is any biosphere-scale structure confluent up to $\rbisim$,
or is confluence strictly small-scale with all large structure race-dominated? If the
latter, confirm that Smith's thermodynamic macro-component and our determinacy
macro-components \emph{cross} rather than nest, and characterise each as the carrier of a
different invariant. \emph{(vi)} Test the candidate that some local confluent island sits at
ecosystem scale against a concrete ecological model, and locate the apertures that bound it.

\section{Conclusion}

We asked where agency lives, and answered with a computational coordinate: at the apertures,
the cuts where nondeterminism is open and observable, which are the races. The route there
was a reframing --- nondeterminism as a resource, a choice in the exact sense, graded by the
axiom of choice and climbing, with its strength, the stalks of a hypercomputation fibration
over the category of computational models and their bisimulation-preserving morphisms. The
mobile process calculi earned their central place by being fundamentally non-confluent, hence
able to speak of races that matter --- two atoms, two sperm, two predators, two buyers --- and
RSpace earned its place by making the choice inside those races visible and runnable. Smith
contributed the move we needed: that the agency boundary is negotiable, that ``agent'' need not
mean ``organism,'' and that one may redraw it --- he, by thermodynamics, around the whole
biosphere. We redraw it instead by determinacy, around the confluent regions, in the manner of
the $\lambda$-calculus; and since selection and predation make Smith's interior anything but
confluent, our boundary and his cross rather than nest. Those confluent regions ---
macro-components --- are small, local islands in a sea of races, and capacity and choice alike
are concentrated at their boundaries, not their interiors.

The deepest claim is the flat ontology, and it is where the two axes most need each other. The
fibration is a chart; the world it charts may be, like the physical and the living world, a
juxtaposition of scales and capacities rather than a tower --- halting oracles beside
$\lambda$-terms, femtosecond chemistry beside minute-scale gene expression, all at the same
cuts. Computation, taught by biology, learns to expect this; biology, taught by computation,
learns what it costs and where it can bite. And the single open hinge --- whether a strong
agent, reaching through an aperture, can crack open the determinism of a weaker one --- is one
confluence lemma, which would tell us, in one stroke and on both axes, whether agency is a fixed
allotment or something agents continually make and unmake in each other.

\section{Coda: hard to see, and the precedent of Bell}\label{sec:coda}

If this picture is right --- if, in the physical world, the surplus power of a high-capacity
agent shows up as nothing more dramatic than the way a low-capacity agent's nondeterministic
interactions get resolved --- then hypercomputation in nature would be extraordinarily hard to
see, and the difficulty would be principled rather than incidental. A confluent region hides
its choices by construction (\S\ref{sec:macro}): the resolution happens but reconverges,
leaving no trace. And even at an aperture, a single resolution looks like nothing --- one
outcome among those the situation allowed, indistinguishable, taken alone, from a coin flip.
Choice-as-resource emits no per-event signal. Whatever evidence there is cannot live in any
single outcome; it must live in the \emph{correlations across} outcomes. The signature of a
resolution stronger than a local choice function is that separated apertures come out
correlated in ways no local, independently-resolving assignment could produce. That is why the
resource is hard to test --- and it tells us exactly where to look.

It is, almost word for word, the lesson of Bell. Bell's theorem shows that no theory in which
each measurement outcome is fixed by local, pre-assigned values --- no \emph{local} choice
function for the outcomes --- can reproduce the correlations quantum mechanics predicts for
separated measurements on an entangled pair \cite{bell}; and the inequalities that mark the
boundary have been violated in the laboratory, culminating in loophole-free experiments
\cite{hensen}. Read in our vocabulary, a quantum measurement is an aperture: a cut whose
outcome the local term does not fix. Bell is then a statement about \emph{correlated choice} ---
that the joint resolution of two separated apertures carries a correlation no local scheduler,
no locally-supplied choice function, can account for. The per-event view sees only randomness
at each wing; the structure is in the correlation, precisely as the testability argument above
predicts. And the correlations have a developed mathematics: the device-independent resource
theory of nonlocality treats these correlations as a graded resource \cite{brunner}, and the
sheaf-theoretic account of contextuality \cite{abramsky} --- logic-and-computer-science
apparatus, already in our bibliography --- characterises exactly the obstruction to a global,
locally-consistent assignment. The bridge between the choice mathematics and the physics already
has a pier on the far bank.

We are not claiming that Bell violations exhibit hypercomputation, or that nonlocal correlations
compute anything uncomputable: they do not, and the resource that nonlocality supplies is not
Turing degree. The precedent is methodological, and in spirit it is decisive. It establishes
that \emph{correlated choice is a real, mathematically characterised, experimentally confirmed
feature of the physical world} --- that a phenomenon which is invisible per-event, visible only
in correlation, and naturally described as the nonlocal resolution of an otherwise-underdetermined
interaction, is not a fantasy of testability but the best-confirmed strange fact in physics.
A research posture that treats choice as a physical resource, and looks for its signature in
correlations across apertures rather than in single events, has therefore already succeeded once,
spectacularly. That is reason to think it could succeed again, for resources further up the
hierarchy than nonlocality reaches.

This is the note we want to leave the reader with, and it points in both directions, as the whole
essay has tried to. Bringing the mathematics of choice --- the axiom of choice and its gradations,
the Weihrauch lattice, the resource theories --- into contact with a physics informed by computation
and biology may be fertile for each. The mathematics gains a fresh set of intuitions and motivations:
choice stops being only a question in the foundations of sets and becomes a calculus of physical and
computational resource, with the strength of a choice principle reading as the altitude it buys in a
hierarchy of capacity, and with empirical stakes. The physics gains a new set of tools: the category
of models and its fibration of hypercomputation, choice-as-resource as the vertical coordinate, the
aperture and the macro-component as a vocabulary for where in a system agency and capacity actually
sit, and a flat ontology in which entities of wildly unequal computational power --- a halting oracle
and a $\lambda$-term --- can be imagined to meet, as so much else in nature meets, at a single cut.

\paragraph{Acknowledgements.}
This essay draws on the RSpace denotation of the rho calculus \cite{qcs,paths} and the GSLT/OSLF
programme \cite{oslf}, and on Smith and Morowitz's reading of the biosphere \cite{smith}; its
conceptual claims are deliberately routed back to structural obligations where they can be settled.

\begin{thebibliography}{99}

\bibitem{rho} L.~Gregory Meredith and Matthias Radestock. A reflective higher-order calculus.
\emph{Electr.\ Notes Theor.\ Comput.\ Sci.}, 141(5):49--67, 2005.

\bibitem{qcs} L.~G. Meredith. \emph{Quoting is Colour-Swap: a model of the rho calculus in the
knotted universe}. Manuscript, 2026.

\bibitem{paths} L.~G. Meredith. \emph{Paths are Subspaces: a cut-triggered polymorphism of RSpace
over path-keys, and its distributive law}. Manuscript, 2026.

\bibitem{knot} L.~G. Meredith. \emph{A Knotted Universe: a new notion of reflective set theories}.
Manuscript, 2026.

\bibitem{oslf} L.~G. Meredith, Christian Wells, et al. Minimal-context transition systems and the
observer-specified logic functor for graph-structured lambda theories. Manuscript, forthcoming.

\bibitem{smith} Eric Smith and Harold J. Morowitz. \emph{The Origin and Nature of Life on Earth:
The Emergence of the Fourth Geosphere}. Cambridge University Press, 2016.

\bibitem{milnerCC} Robin Milner. \emph{Communication and Concurrency}. Prentice Hall, 1989.

\bibitem{leifermilner} James J. Leifer and Robin Milner. Deriving bisimulation congruences for
reactive systems. In \emph{CONCUR 2000}, LNCS 1877, pages 243--258. Springer, 2000.

\bibitem{sewell} Peter Sewell. From rewrite rules to bisimulation congruences. \emph{Theoret.\
Comput.\ Sci.}, 274(1--2):183--230, 2002.

\bibitem{turing1939} Alan M. Turing. Systems of logic based on ordinals. \emph{Proc.\ London Math.\
Soc.}, s2-45(1):161--228, 1939.

\bibitem{soare} Robert I. Soare. \emph{Turing Computability: Theory and Applications}. Springer,
2016.

\bibitem{copeland} B. Jack Copeland. Hypercomputation. \emph{Minds and Machines}, 12(4):461--502, 2002.

\bibitem{weihrauch} Vasco Brattka, Guido Gherardi, and Arno Pauly. Weihrauch complexity in computable
analysis. In V. Brattka and P. Hertling, editors, \emph{Handbook of Computability and Complexity in
Analysis}, pages 367--417. Springer, 2021.

\bibitem{rutten} J.~J.~M.~M. Rutten. Universal coalgebra: a theory of systems. \emph{Theoret.\
Comput.\ Sci.}, 249(1):3--80, 2000.

\bibitem{aczel} Peter Aczel. \emph{Non-Well-Founded Sets}, CSLI Lecture Notes 14. CSLI Publications,
1988.

\bibitem{bell} John S. Bell. On the Einstein Podolsky Rosen paradox. \emph{Physics Physique
Fizika}, 1(3):195--200, 1964.

\bibitem{hensen} B. Hensen et al. Loophole-free Bell inequality violation using electron spins
separated by 1.3 kilometres. \emph{Nature}, 526:682--686, 2015.

\bibitem{brunner} Nicolas Brunner, Daniel Cavalcanti, Stefano Pironio, Valerio Scarani, and
Stephanie Wehner. Bell nonlocality. \emph{Rev.\ Mod.\ Phys.}, 86:419--478, 2014.

\bibitem{abramsky} Samson Abramsky and Adam Brandenburger. The sheaf-theoretic structure of
non-locality and contextuality. \emph{New J.\ Phys.}, 13:113036, 2011.

\end{thebibliography}

\end{document}
